% GENERATED FILE. Edit canonical lessons, then run npm run generate:books.
% canonical-source-hash: fnv1a64:695c3e1e74153185
% canonical-lessons: TA-C257-veru, TA-C257-accariyappatu, TA-C257-vetkappatu, TA-C257-poramaippatu, TA-C257-muttamitu

\chapter{To hate, To be surprised, To feel shy, To be jealous, To kiss}
\label{ch:ta-to-hate-to-be-surprised-to-feel-shy-to-be-jealous-to-kiss}
\hlchaptermodality{257}

\begin{chapteropening}
\textbf{By the end of this chapter:} \emph{I can say to hate, to be surprised, to feel shy, to be jealous and to kiss.}

Complete the last lesson of chapter 257: I can say to hate, to be surprised, to feel shy, to be jealous and to kiss.
\end{chapteropening}

\section[\texorpdfstring{\ta{வெறு} (veṟu)}{வெறு (veṟu)}]{\ta{வெறு} (veṟu) — to hate}

\label{lesson:TA-C257-veru}



\begin{quote}
Before the new one: say the Tamil for to erase, then the Tamil for to print.
\end{quote}

\subsection*{You'll want to know: \ta{வெறு}}
\textbf{\ta{வெறு}} — \emph{veṟu} — \textquotedblleft{}to hate\textquotedblright{}.

\begin{grammarlens}[title={What you've built}]
One more word for this chapter.
\end{grammarlens}

\subsection*{Your turn}
\begin{itemize}
\raggedright
  \item \emph{Say it:} \emph{veṟu}
  \item \emph{Say it:} \emph{veṟu}, once more
  \item \emph{Say it:} say \emph{veṟu}
  \item \emph{Recall:} say the Tamil for to erase, then the Tamil for to print, then say \emph{veṟu} again
\end{itemize}

\begin{tcolorbox}[breakable,colback=teal!4,colframe=teal!35!black,title={Before you move on}]
What does \ta{வெறு} mean? (\textquotedblleft{}To hate\textquotedblright{}.) Say it once more.
\end{tcolorbox}
\section[\texorpdfstring{\ta{ஆச்சரியப்படு} (āccariyappaṭu)}{ஆச்சரியப்படு (āccariyappaṭu)}]{\ta{ஆச்சரியப்படு} (āccariyappaṭu) — to be surprised}

\label{lesson:TA-C257-accariyappatu}



\begin{quote}
Before the new one: say the Tamil for to print, then the Tamil for to hate.
\end{quote}

\subsection*{You'll want to know: \ta{ஆச்சரியப்படு}}
\textbf{\ta{ஆச்சரியப்படு}} — \emph{āccariyappaṭu} — \textquotedblleft{}to be surprised\textquotedblright{}.

\begin{grammarlens}[title={What you've built}]
One more word for this chapter.
\end{grammarlens}

\subsection*{Your turn}
\begin{itemize}
\raggedright
  \item \emph{Say it:} \emph{āccariyappaṭu}
  \item \emph{Say it:} \emph{āccariyappaṭu}, once more
  \item \emph{Say it:} say \emph{āccariyappaṭu}
  \item \emph{Recall:} say the Tamil for to print, then the Tamil for to hate, then say \emph{āccariyappaṭu} again
\end{itemize}

\begin{tcolorbox}[breakable,colback=teal!4,colframe=teal!35!black,title={Before you move on}]
What does \ta{ஆச்சரியப்படு} mean? (\textquotedblleft{}To be surprised\textquotedblright{}.) Say it once more.
\end{tcolorbox}
\section[\texorpdfstring{\ta{வெட்கப்படு} (veṭkappaṭu)}{வெட்கப்படு (veṭkappaṭu)}]{\ta{வெட்கப்படு} (veṭkappaṭu) — to feel shy, to be embarrassed}

\label{lesson:TA-C257-vetkappatu}



\begin{quote}
Before the new one: say the Tamil for to hate, then the Tamil for to be surprised.
\end{quote}

\subsection*{You'll want to know: \ta{வெட்கப்படு}}
\textbf{\ta{வெட்கப்படு}} — \emph{veṭkappaṭu} — \textquotedblleft{}to feel shy, to be embarrassed\textquotedblright{}.

\begin{grammarlens}[title={What you've built}]
One more word for this chapter.
\end{grammarlens}

\subsection*{Your turn}
\begin{itemize}
\raggedright
  \item \emph{Say it:} \emph{veṭkappaṭu}
  \item \emph{Say it:} \emph{veṭkappaṭu}, once more
  \item \emph{Say it:} say \emph{veṭkappaṭu}
  \item \emph{Recall:} say the Tamil for to hate, then the Tamil for to be surprised, then say \emph{veṭkappaṭu} again
\end{itemize}

\begin{tcolorbox}[breakable,colback=teal!4,colframe=teal!35!black,title={Before you move on}]
What does \ta{வெட்கப்படு} mean? (\textquotedblleft{}To feel shy, to be embarrassed\textquotedblright{}.) Say it once more.
\end{tcolorbox}
\section[\texorpdfstring{\ta{பொறாமைப்படு} (poṟāmaippaṭu)}{பொறாமைப்படு (poṟāmaippaṭu)}]{\ta{பொறாமைப்படு} (poṟāmaippaṭu) — to be jealous}

\label{lesson:TA-C257-poramaippatu}



\begin{quote}
Before the new one: say the Tamil for to be surprised, then the Tamil for to feel shy.
\end{quote}

\subsection*{You'll want to know: \ta{பொறாமைப்படு}}
\textbf{\ta{பொறாமைப்படு}} — \emph{poṟāmaippaṭu} — \textquotedblleft{}to be jealous\textquotedblright{}.

\begin{grammarlens}[title={What you've built}]
One more word for this chapter.
\end{grammarlens}

\subsection*{Your turn}
\begin{itemize}
\raggedright
  \item \emph{Say it:} \emph{poṟāmaippaṭu}
  \item \emph{Say it:} \emph{poṟāmaippaṭu}, once more
  \item \emph{Say it:} say \emph{poṟāmaippaṭu}
  \item \emph{Recall:} say the Tamil for to be surprised, then the Tamil for to feel shy, then say \emph{poṟāmaippaṭu} again
\end{itemize}

\begin{tcolorbox}[breakable,colback=teal!4,colframe=teal!35!black,title={Before you move on}]
What does \ta{பொறாமைப்படு} mean? (\textquotedblleft{}To be jealous\textquotedblright{}.) Say it once more.
\end{tcolorbox}
\section[\texorpdfstring{\ta{முத்தமிடு} (muttamiṭu)}{முத்தமிடு (muttamiṭu)}]{\ta{முத்தமிடு} (muttamiṭu) — to kiss}

\label{lesson:TA-C257-muttamitu}



\begin{quote}
Before the new one: say the Tamil for to feel shy, then the Tamil for to be jealous.
\end{quote}

\subsection*{You'll want to know: \ta{முத்தமிடு}}
\textbf{\ta{முத்தமிடு}} — \emph{muttamiṭu} — \textquotedblleft{}to kiss\textquotedblright{}.

\begin{grammarlens}[title={What you've built}]
One more word for this chapter.
\end{grammarlens}

\subsection*{Your turn}
\begin{itemize}
\raggedright
  \item \emph{Say it:} \emph{muttamiṭu}
  \item \emph{Say it:} \emph{muttamiṭu}, once more
  \item \emph{Say it:} say \emph{muttamiṭu}
  \item \emph{Recall:} say the Tamil for to feel shy, then the Tamil for to be jealous, then say \emph{muttamiṭu} again
\end{itemize}

\begin{tcolorbox}[breakable,colback=teal!4,colframe=teal!35!black,title={Before you move on}]
What does \ta{முத்தமிடு} mean? (\textquotedblleft{}To kiss\textquotedblright{}.) Say it once more.
\end{tcolorbox}
